% ============================================================
%  参考文献
% ============================================================
\addcontentsline{toc}{section}{参考文献}

\begin{thebibliography}{AAEKO}

\bibitem[AEK]{AEK}
D.~Auroux, A.~Efimov, L.~Katzarkov,
\emph{Lagrangian Floer theory for trivalent graphs and homological mirror symmetry for curves},
preprint, arXiv:2107.01981.

\bibitem[AAEKO]{AAEKO}
M.~Abouzaid, D.~Auroux, A.~Efimov, L.~Katzarkov, D.~Orlov,
\emph{Homological mirror symmetry for punctured spheres},
J.~Amer. Math. Soc. \textbf{26} (2013), 1051--1083.

\bibitem[AAK]{AAK}
M.~Abouzaid, D.~Auroux, L.~Katzarkov,
\emph{Homological mirror symmetry for hypersurfaces in $(\C^*)^n$},
preprint.

\bibitem[HV]{HV}
K.~Hori, C.~Vafa,
\emph{Mirror symmetry},
preprint, hep-th/0002222.

\bibitem[Kon]{Kon}
M.~Kontsevich,
\emph{Homological algebra of mirror symmetry},
Proceedings of the International Congress of Mathematicians (Z\"urich, 1994), Birkh\"auser, 1995, 120--139.

\bibitem[Kn]{Kn}
H.~Kn\"orrer,
\emph{Cohen--Macaulay modules on hypersurface singularities I},
Invent. Math. \textbf{88} (1987), 153--164.

\bibitem[Or]{Or}
D.~Orlov,
\emph{Derived categories of coherent sheaves and triangulated categories of singularities},
in: Algebra, Arithmetic and Geometry, Progr. Math. \textbf{270}, Birkh\"auser, 2009, 503--531.

\bibitem[PZ]{PZ}
A.~Polishchuk, E.~Zaslow,
\emph{Categorical mirror symmetry: the elliptic curve},
Adv. Theor. Math. Phys. \textbf{2} (1998), 443--470; arXiv:math/9801119.

\bibitem[Sei]{Sei}
P.~Seidel,
\emph{Fukaya Categories and Picard--Lefschetz Theory},
Zurich Lectures in Advanced Mathematics, European Mathematical Society, 2008.

\bibitem[AS]{AS}
M.~Abouzaid, P.~Seidel,
\emph{An open string analogue of Viterbo functoriality},
Geom. Topol. \textbf{14} (2010), 627--718.

\bibitem[Aur1]{Aur1}
D.~Auroux,
\emph{Mirrors of curves and their Fukaya categories},
迈阿密，2020 年 1 月 27--29 日三场报告的手写笔记（17 页）。本册第 1--3 讲的底本。

\bibitem[Aur2]{Aur2}
D.~Auroux,
\emph{An invitation to homological mirror symmetry},
IMSA 成立大会，迈阿密，2019 年 9 月 6 日，幻灯片 14 页。本册附录 A 的底本。

\end{thebibliography}

\noindent\textbf{译者附识.} 原稿在正文中常按姓氏提及他人的工作（Hanlon 的学位论文、J. Hicks 的热带拉格朗日、N. Jeffs 的 Kn\"orrer 周期性、C. Cannizzo 的学位论文、H. Lee、Gammage--Shende、Lekili--Polishchuk、Ganatra 等）。这些工作原稿未给出可核验的完整书目信息，本表不臆补条目；读者可按作者与主题在原稿所指处检索。
